\documentclass[12pt]{article}
\usepackage[margin=1in]{geometry}
\usepackage{amsmath,amsthm,amssymb,scrextend}
\setlength{\parindent}{0pt}
\begin{document}
\newcommand{\ds}{\displaystyle}
Let $T$ be the waiting time until the first text, and let $N(t)$ be the number of texts in the first $t$ hours

$$N(t)\sim \operatorname{Poisson}(\lambda t)$$

$$T>t\iff N(t)=0$$

$$P(T>t)=P(N(t)=0)=e^{-\lambda t}$$

$$p=P(1<T<2)=P(T>1)-P(T>2)=e^{-\lambda}-e^{-2\lambda}$$
\bigskip\hrule\bigskip

Let $\ds x=e^\lambda$

$$px^2-x+1=0$$

$$x=\frac{1\pm\sqrt{1-4p}}{2p}$$

$\because 0<p<\frac14$

$$1-4p>0$$

so there are two distinct real roots.

$$\frac{1+\sqrt{1-4p}}{2p}>1,\qquad \frac{1-\sqrt{1-4p}}{2p}=\frac{2}{1+\sqrt{1-4p}}>1$$

$$\therefore x>1$$
$$\therefore \lambda=\ln x>0$$

There are two positive values of $\lambda$

\bigskip\hrule\bigskip

$e^{\lambda_1}$ and $e^{\lambda_2}$ are roots of $px^2-x+1=0$

$$e^{\lambda_1}e^{\lambda_2}=\frac1p\implies \lambda_1+\lambda_2=\ln\frac1p$$

The two phones together have rate of $\lambda_1+\lambda_2$

$$P(1<T<2)=e^{-(\lambda_1+\lambda_2)}-e^{-2(\lambda_1+\lambda_2)}$$

$$\because e^{\lambda_1+\lambda_2}=\frac1p$$
$$\therefore \boxed{P(1<T<2)=p-p^2}$$

\bigskip
\hrule
\begin{center}
    This solution was submitted by \textbf{Jon}\\[4pt]
    Published and refined by \textbf{Jon}\\[8pt]
    \textit{Last updated: June 25, 2026}
\end{center}

\end{document}
